\documentclass[11pt]{article}
\usepackage[margin=1in]{geometry}
\usepackage{amsmath,amssymb}
\usepackage{listings}
\usepackage{hyperref}
\usepackage{longtable}
\usepackage{booktabs}
\lstset{basicstyle=\ttfamily\footnotesize,breaklines=true,columns=flexible}

\title{Independent Replication of ``Quantum Algorithm Implementations for Beginners''\\
{\normalsize arXiv:1804.03719v3 --- Abhijith J.\ et al., LANL, 2018/2022}}
\author{Independent replication (QC-200 wave, 2026-07-05)}
\date{2026-07-05}

\begin{document}
\maketitle

\section*{Verdict (one line)}
\textbf{SPOT-CHECK.} All three independently chosen algorithms
(Bernstein--Vazirani, Grover, Quantum Phase Estimation) reproduce
their expected textbook outputs to statevector precision using a
real Qiskit~2.5.0 statevector simulation, matching the paper's
qualitative description in Sections~2, 3, and~4.3. The paper is a
20-algorithm pedagogical survey with no single headline number, so
this is a legitimate SPOT-CHECK (not full REPLICATED, which would
require reproducing every algorithm plus every hardware-vs.-sim
histogram in the paper).

\section{Paper summary}
\label{sec:summary}
Abhijith~J.\ et~al.\ (LANL, LA-UR-20-22353) is a 100+ page
pedagogical review that presents 20 quantum algorithms
with reference IBM Qiskit implementations, and for each
algorithm compares outputs on the noiseless Qiskit simulator vs
the noisy IBM Q~4/Q~5 hardware available in the 2018--2020 window.
Algorithms covered include Grover's search, Bernstein--Vazirani
(BV), Simon's algorithm, Deutsch--Jozsa, Quantum Phase Estimation
(QPE), Shor's factoring, HHL for linear systems, QAOA for
MaxCut, VQE for small Hamiltonians, and several more.

The paper has no single ``headline number''; it is a survey.
Its verifiable content is therefore per-algorithm: each of the 20
sections gives an expected output shape (which measurement
outcome should dominate) and a hardware histogram. A faithful
independent reproduction picks a small subset of algorithms and
verifies the paper's expected \emph{simulator-level} outcomes
end-to-end.

\section{Claims table}
\label{sec:claims}
\begin{longtable}{@{}p{2.6cm}p{6.4cm}cc@{}}
\toprule
Claim id & Statement & Testable? & Tested? \\
\midrule
C1 (BV) & For a linear oracle $f(x)=s\cdot x \pmod 2$ with hidden
$s\in\{0,1\}^n$, one query of the BV circuit recovers $s$ with
probability~1 on a noiseless simulator (Sec.~3). & Yes & Yes \\
C2 (Grover) & For $N=2^n$ and $M=1$ marked item, after
$k^*=\lfloor(\pi/4)\sqrt{N/M}\rceil$ Grover iterations, the marked
state has probability $\sin^2((2k^*+1)\theta)$ with
$\theta=\arcsin\sqrt{M/N}$. For $n=3$, $M=1$: $k^*=2$ and
$P\approx0.9453$ (Sec.~2). & Yes & Yes \\
C3 (QPE) & For a 1-qubit unitary with eigenphase $\varphi=k/2^t$
exactly representable in $t$ counting qubits, QPE returns the
integer $k$ with probability~1 (Sec.~4.3). & Yes & Yes \\
C4 (Survey scope) & The paper presents Qiskit-compatible
reference implementations for 20 algorithms. & No (assertion of
existence, not a numerical claim). & Not applicable \\
C5 (Hardware vs sim) & On the IBM Q4/Q5 hardware of 2018--2020,
noiseless-sim outcomes differ measurably from hardware
histograms. & Yes, but requires IBM hardware access that is now
retired for those chips. & No (blocked). \\
\bottomrule
\end{longtable}

\section{Method}
\label{sec:method}
\subsection{Environment}
\begin{itemize}
\item Host: Darwin 25.3.0 (CherryRd), CPU only.
\item Python 3.12.13 in \texttt{venv/}.
\item Qiskit 2.5.0, qiskit-aer (via \texttt{pip install qiskit qiskit-aer numpy}), NumPy 2.5.1.
\item Statevector via \texttt{qiskit.quantum\_info.Statevector.from\_instruction} --- no shot noise, no sampler.
\end{itemize}

\subsection{Bernstein--Vazirani (C1)}
Hidden string $s = 1011$ (MSB-first). Circuit uses $n{=}4$ input
qubits and 1 ancilla in $|-\rangle$. Oracle implemented by
$\text{CX}(x_i \to a)$ for each $s_i=1$. Full script at
\texttt{report/evidence/bv.py}. Runs in $\ll1$~s.

\subsection{Grover (C2)}
$n{=}3$, marked state $|101\rangle$ (decimal~5). Oracle:
X-flip qubits where $m$ has a 0, then $H \cdot \text{MCX} \cdot H$
on the top qubit (implements a multi-controlled $Z$), then undo
the flips. Diffuser: standard $H X (H\text{MCX}H) X H$
construction. $k^*=2$ iterations. Full script at
\texttt{report/evidence/grover.py}.

\subsection{QPE (C3)}
Target $U=\text{diag}(1,e^{2\pi i\varphi})$ with $\varphi=1/8$,
so $2^t\varphi = 16\cdot1/8 = 2$ is an integer for $t=4$.
Eigenstate register prepared as $|1\rangle$. Counting register:
$t=4$ qubits, controlled-$U^{2^j}$ from qubit $j$, then inverse
QFT (explicit implementation with swaps + graded controlled
phases). Full script at \texttt{report/evidence/qpe.py}.

\section{Results vs paper}
\label{sec:results}
\begin{longtable}{@{}p{2.5cm}p{4.5cm}p{4.5cm}c@{}}
\toprule
Claim & Paper prediction & Our simulation & Match? \\
\midrule
C1 (BV, $s{=}1011$) & Recover $s$ with $P=1$ in 1 query & Recover
$1011$ with $P=0.9999\underline{9999999998}6$ (i.e.\ 1 up to
double precision) & \textbf{YES} \\
C2 (Grover $N{=}8$) & After $k^*=2$: $P(\text{marked})\approx
0.9453$ & $P(|101\rangle)=0.9453124999999959$ (analytic
$=0.9453124999999999$) & \textbf{YES} \\
C3 (QPE $\varphi{=}1/8$) & $|0010\rangle$ with $P=1$ & $|0010\rangle$
with $P=0.9999\underline{9999999998}7$; all other basis states
have probability $\le 3\times10^{-31}$ & \textbf{YES} \\
\bottomrule
\end{longtable}

Full JSON outputs (with the entire probability distribution over
all $2^n$ basis states in each case) are saved to:
\begin{itemize}
\item \texttt{report/evidence/bv\_result.json}
\item \texttt{report/evidence/grover\_result.json}
\item \texttt{report/evidence/qpe\_result.json}
\end{itemize}

\section{Verdict + justification}
\label{sec:verdict}
\textbf{SPOT-CHECK.} All three independently chosen algorithms
match their textbook / paper-stated simulator predictions to
statevector precision (double). The paper is a survey with no
single headline number, so full ``REPLICATED'' would require
reproducing all 20 algorithms plus each hardware histogram; that
was out of scope for this timed replication. What is claimed here
is: the paper's Qiskit-style pseudocode for these three
canonical algorithms is faithful to the underlying quantum theory,
and independent implementations (written from the paper's
descriptions without copying code) reproduce the promised
noiseless outputs exactly.

\section{Open Questions}
\label{sec:openq}
See also \texttt{report/open\_questions.json} for the
machine-readable form.

\begin{enumerate}
\item[Q1] The paper reports per-algorithm hardware histograms
from IBM Q4/Q5 chips (2018--2020). Those chips have since been
retired. How would today's ibmq\_qasm equivalents (Falcon/Eagle)
change the hardware-vs-sim gaps the paper documented, and can any
of the ``worked on hardware'' verdicts be reversed by
present-day noise floors?
\item[Q2] The QPE section uses $t$ counting qubits with the
implicit assumption that $2^t\varphi$ is an integer for the
paper's demo. What is the probability of the second- and
third-most-likely outcomes when $\varphi$ is \emph{not}
representable exactly in $t$ bits? A parameter sweep would
quantify how the paper's ``exact'' demo degrades to the generic
$1/(k^2\pi^2)$-tailed distribution, which the paper does not
plot.
\item[Q3] The Grover section quotes success probability $>1/2$
as ``the guarantee'' but the actual $k^*=\lfloor(\pi/4)\sqrt{N/M}\rceil$
rounding can either over- or under-shoot the optimal angle.
Which specific $(N,M)$ pairs land closer to $1/2$ than to $1$?
This has practical consequence for the paper's IBM-hardware
runs but is not tabulated in the paper.
\item[Q4] The paper's Bernstein--Vazirani implementation uses
$n{+}1$ qubits and 1 oracle call, achieving $P=1$ on a
simulator. The paper does not quantify the crossover point at
which BV becomes hardware-\emph{worse} than classical bit-by-bit
probing under a given depolarizing-noise rate. A depolarizing
sweep across $n=2\ldots16$ would answer whether the paper's
``quantum advantage'' argument survives on today's noisiest
public chips.
\item[Q5] The paper's own extraction (Marker on the same PDF)
regularly bleeds the first author's surname ``ABHIJITH'' into
the title. This is a corpus-hygiene issue that likely affects
downstream tools using the paper's Marker parse as ground
truth. What fraction of the paper's cited references would be
mis-attributed if a metadata pipeline used the Marker-parsed
title as the paper's true title?
\end{enumerate}

\end{document}
